\documentclass[12pt,a4paper]{article}
\usepackage{amsmath,amssymb,amsthm}
\usepackage{geometry}
\usepackage{hyperref}
\usepackage{microtype}
\usepackage{booktabs}
\usepackage{array}
\geometry{margin=2.5cm}

\newtheorem{theorem}{Theorem}[section]
\newtheorem{lemma}[theorem]{Lemma}
\newtheorem{corollary}[theorem]{Corollary}
\newtheorem{definition}[theorem]{Definition}
\newtheorem{remark}[theorem]{Remark}
\newtheorem{observation}[theorem]{Observation}

\title{The Cascade Series --- Part II\\[0.5em]
\large Quantum Mechanics from the Cascade:\\
Effective Theory of a 4-Dimensional Observer\\
in the Sphere-Area Geometry}
\author{RTAC}
\date{March 2026}

\begin{document}
\maketitle

\begin{abstract}
The cascade series tests one hypothesis: the infinite-dimensional unit ball, descended to
four dimensions, is indistinguishable from our universe. The companion paper~\cite{paperI}
derived a geometric invariant $I = 1.0996\times 10^{-120}$ from the sphere-area cascade
using only orthogonality as input. Here we show that a 4-dimensional observer embedded in
the cascade geometry recovers the structural framework of quantum mechanics without
additional postulates. The argument proceeds in five steps. (1)~Projection of the
high-dimensional unit-sphere measure onto a 4D subspace produces Gaussian statistics via
the central limit theorem inherited from concentration of measure. (2)~The cascade's
orthogonal slicing structure induces orthogonal measurement bases in the projected theory.
(3)~Concentration of measure on $S^{d-1}$ for large $d$ recovers the Born rule as the
unique consistent probability assignment, reproducing Gleason's theorem geometrically.
(4)~The precession angle between consecutive slicing axes is $\alpha=\pi/2$, forced by the
same orthogonality condition that forces $\sqrt{\pi}$ in the slicing recurrence
(Theorem~\ref{thm:precession}). The proof distinguishes two applications of the same
geometric fact: the slicing axis is perpendicular to the equatorial hyperplane (forcing
$\Gamma(\tfrac{1}{2})=\sqrt{\pi}$ in the integral), and the subsequent slicing axis is
perpendicular to the already-integrated direction (forcing $\alpha=\pi/2$ for the
precession). These are the same orthogonality axiom applied at successive steps; the
cascade's axiom forces both simultaneously. (5)~The identification of time with the
cascade's orthogonal slicing direction yields a discrete propagator whose continuum limit
is a Schr\"{o}dinger-type evolution equation, with complexification arising from the forced
precession.

Part~0 establishes that sphere areas are the unique independent cascade quantities
(Corollary~3.2). This paper proves the compactification radius
$R_{\rm eff}=1/\sqrt{d+3}$ from the Beta function and develops the full geometric
interpretation: each slicing step sends one direction's effective radius to
$R_{\rm eff}$, giving concrete geometric
content to the arrow of time as asymptotic suppression rather than sharp elimination. The
cascade's two natural length scales---the compactification radius and the Gaussian width of
the slicing integrand---coincide asymptotically, producing a second universal cascade constant
$\operatorname{erf}(1/\sqrt{2}) = 0.68269$. The boundary dominance
$\Omega_{d-1}/V_d=d$ (Theorem~3.1 of~\cite{paperI}) confirms that the observer's state
space is the unit sphere because the cascade's content is its boundaries. The cascade fixes
the arena (which Hilbert space), the measure (why Gaussian), and the normalisation constant
($\sqrt{\pi}$) that standard quantum mechanics postulates. The value of $\hbar$ is not
derived; it enters as the ratio of the forced precession scale $\pi/2$ to the physical time
increment, and therefore encodes only the unit matching between geometric and physical time.
No physical law or free geometric parameter beyond the observer's dimensionality enters the
derivation.
\end{abstract}

\tableofcontents
\newpage

\section{Introduction and Summary of the Foundation}

In~\cite{paperI}, the sphere-area cascade was derived from a single axiom: orthogonality.
Starting from the infinite-dimensional unit ball---which has zero volume, zero surface area,
and no interior---the cascade compresses orthogonal directions one at a time, descending
from $d=\infty$. The unique dimension-independent constant in the slicing recurrence is
$\Gamma(\tfrac{1}{2})=\sqrt{\pi}$, forced by the quarter-turn integral $B(\tfrac{1}{2},\cdot)$.
This constant enters in two roles (multiplicative factor and bare decay rate), producing two
distinguished dimensions $d_1=19$ and $d_2=217$, and a cascade invariant
$I=\Omega_{19}\times\Omega_{217}\times 198/217=1.0996\times 10^{-120}$.

Corollary~3.2 of~\cite{paperI} establishes that sphere areas are the unique independent
cascade quantities: every other cascade object is derived from sphere areas, and any
identification mapping cascade content to energy must do so with a universal constant.
Section~\ref{sec:compactification} of this paper develops the full compactification
interpretation: each slicing step sends one direction's effective radius to
$R_{\rm eff}(d)=1/\sqrt{d+3}$, and the boundary-to-volume ratio satisfies
$\Omega_{d-1}/V_d=d$ for all $d\geq 1$. The cascade's content is its boundaries; sphere
areas, not volumes, are the primary objects.

The present paper asks: what does the cascade geometry look like from inside? Specifically,
if an observer is restricted to a 4-dimensional cross-section of the full cascade structure,
what measurement theory does that observer use? We show the answer is quantum mechanics.

The logical structure is:
\begin{itemize}
\item \emph{Foundation~\cite{paperI}}: Orthogonality $\to$ $\sqrt{\pi}$ $\to$ slicing
recurrence $\to$ cascade $\to$ tower $(d_0,d_1,d_2)$ $\to$ $I\approx 10^{-120}$. Each step
is an asymptotic compactification with $R_{\rm eff}=1/\sqrt{d+3}$.
\item \emph{This paper}: Cascade geometry + 4D observer $\to$ Hilbert space structure +
Gaussian measure + Born rule + unitary evolution $=$ QM.
\end{itemize}

The derivation does not explain why we observe four dimensions. That is an empirical input,
just as in~\cite{paperI}. What we show is that, given a 4D perspective, the cascade's
internal structure reproduces the quantum formalism.

\section{The Observer's Arena}

\subsection{State space is the unit sphere}

Quantum states are unit vectors in a Hilbert space $\mathcal{H}$. The set of pure states is
the unit sphere $S(\mathcal{H})$. In finite dimension $d$, this is $S^{d-1}\subset\mathbb{R}^d$
(or $S^{2d-1}\subset\mathbb{C}^d$ after complexification). The cascade geometry provides
exactly this arena: the unit sphere $S^{d-1}$ at each level $d$ of the descent. A 4D
observer---one who can access only 4 orthogonal directions---sees the projection of
$S^{d-1}$ onto a 4-dimensional subspace. This projected space is the observer's effective
state space.

\subsection{Why the cascade fixes the arena}

Standard quantum mechanics postulates a Hilbert space but does not derive it. The cascade
provides the derivation: the arena is the unit sphere because the cascade is a sequence of
volume-preserving orthogonal compressions of the unit ball. Every state accessible to any
observer at any dimension lies on a unit sphere. The constraint $|\psi|^2=1$, which in QM
is the normalisation postulate, is here a theorem: the cascade operates on $B^d$, whose
boundary is $S^{d-1}$.

\subsection{Boundary dominance and the primacy of the sphere}

The arena is the unit sphere rather than the unit ball for a reason that~\cite{paperI},
Section~3 makes precise.

\begin{theorem}[Boundary dominance; Theorem~3.1 of~\cite{paperI}]
$\Omega_{d-1}/V_d=d$ for all $d\geq 1$.
\end{theorem}

At $d=4$: $\Omega_3/V_4=4$, so the boundary $S^3$ carries a fraction $d/(d+1)=4/5=80\%$
of the content of $B^4$. At the cascade dimensions $d_1=19$ and $d_2=217$, the fraction
on the boundary is $19/20=95\%$ and $217/218\approx 99.5\%$ respectively. The cascade's
content is its boundaries, increasingly so at each step. The observer's state space is the
sphere because the sphere carries nearly all the geometric content available to the observer.
This is not an axiom; it is a consequence of the Gamma function identity
$\Omega_{d-1}/V_d = 2\pi^{d/2}/[\Gamma(d/2)] \div \pi^{d/2}/[\Gamma(d/2+1)] = d$.

\begin{remark}[Layered boundary dominance]
The cascade stacks boundary-dominant steps: the $d_2$-sphere's boundary encloses the
$d_1$-sphere's boundary, which encloses the $d_0$-sphere's boundary, which encloses the
observer's $S^3$. At each step, the boundary $S^{d-1}$ carries a fraction $d/(d+1)$ of the
content. The cascade is a nested sequence of shells, each carrying nearly all the geometric
content of its level, with the observer on the innermost shell.
\end{remark}

\section{Gaussian Statistics from Projection}

\subsection{Concentration of the slicing integrand}

The slicing integrand $f_d(x) = (1-x^2)^{d/2}$ that defines the cascade recurrence
is itself approximately Gaussian. This is the concrete mechanism underlying the
projection theorem that follows.

\begin{lemma}[Gaussian concentration of the integrand]\label{lem:gaussian}
The slicing integrand $f_d(x) = (1-x^2)^{d/2}$ satisfies:
\begin{enumerate}
\item[(a)] $f_d(x) = \exp(-dx^2/2 + O(dx^4))$, approximating a Gaussian with standard
deviation $\sigma(d) = 1/\sqrt{d}$;
\item[(b)] the half-maximum width is $\sqrt{2\ln 2}/\sqrt{d}$;
\item[(c)] as $d\to\infty$, $f_d\to\delta(x)$.
\end{enumerate}
\end{lemma}
\begin{proof}
Expanding: $\ln(1-x^2) = -x^2 - x^4/2 - x^6/3 - \cdots$ for $|x|<1$. Therefore
$(d/2)\ln(1-x^2) = -dx^2/2 + O(dx^4)$. The half-maximum occurs at
$x_{1/2} = \sqrt{1-2^{-2/d}} \approx \sqrt{2\ln 2}/\sqrt{d}$.
As $d\to\infty$, $x_{1/2}\to 0$ and $\int f_d\,dx \approx \sqrt{2\pi/d}\to 0$,
confirming convergence to $\delta(x)$.
\end{proof}

This lemma is the elementary origin of all Gaussian structure in the cascade. The
projection theorem (below) elevates it to a distributional statement on the unit
sphere; the compactification results (Section~\ref{sec:compactification}) use it to
define the effective radius of each temporal step.

\subsection{The projection theorem}

\begin{theorem}[Gaussian emergence]\label{thm:gaussian}
Let $x$ be uniformly distributed on $S^{d-1}(\sqrt{d})\subset\mathbb{R}^d$. Let
$\pi_k:\mathbb{R}^d\to\mathbb{R}^k$ be projection onto any $k$-dimensional subspace, with
$k$ fixed and $d\to\infty$. Then $\pi_k(x)$ converges in distribution to
$\mathcal{N}(0,I_k)$, the standard $k$-dimensional Gaussian.
\end{theorem}

This is a consequence of the Poincar\'{e} limit theorem (see~\cite{ledoux,milman}). The
rate of convergence is $O(1/d)$, so for the cascade dimensions $d_1=19$, $d_2=217$, the
approximation is already excellent for $k=4$.

\subsection{The cascade's Gaussian is the QM wavefunction}

The Gaussian that emerges from projection is structurally identical to the ground-state
wavefunction of the quantum harmonic oscillator:
$\psi_0(x)=\pi^{-1/4}\exp(-x^2/2)$. The normalisation constant $\pi^{-1/4}$ involves
exactly the $\sqrt{\pi}$ from the cascade. In standard QM this constant is obtained by
requiring $\int|\psi|^2dx=1$ with the Gaussian integral $\int\exp(-x^2)dx=\sqrt{\pi}$.
In the cascade, $\sqrt{\pi}$ is not a normalisation convention: it is the unique
dimension-independent constant forced by orthogonal compression (Theorem~3.1
of~\cite{paperI}). The cascade derives what QM postulates.

\subsection{Why Gaussians are universal}

The cascade provides a structural explanation for why Gaussians dominate quantum physics.
Any 4D observer embedded in a high-dimensional unit-sphere geometry necessarily sees
Gaussian statistics, because projection from high to low dimensions produces Gaussians.
This is a geometric statement: the shape of the wavefunction is fixed by the shape of
the arena.

\section{Orthogonal Measurement Bases}

\subsection{Slicing induces measurement}

In the cascade, each step selects an orthogonal direction and integrates over it. From the
perspective of a 4D observer, this operation has a natural interpretation: choosing a
direction along which to slice is choosing an observable to measure. The act of slicing---
projecting the $(d+1)$-ball onto the equatorial $d$-ball---discards information about the
perpendicular coordinate. This is measurement: extracting a definite value along one axis
at the cost of information about the complementary axis.

The compactification results of Section~\ref{sec:compactification} sharpen this
interpretation. Each slicing step does not annihilate the integrated-out direction; it
suppresses it to an effective radius $R_{\rm eff}(d)=1/\sqrt{d+3}$
(Theorem~\ref{thm:reff}). The weight function $(1-x^2)^{d/2}$ remains positive on $(-1,1)$
for all finite $d$. Measurement in the cascade is asymptotic suppression, not sharp
projection: the complementary coordinate is exponentially suppressed but never exactly zero.
This is a geometric realisation of the fact that quantum measurement does not destroy the
unmeasured degree of freedom---it renders it inaccessible at scale $R_{\rm eff}$.

\begin{theorem}[Orthogonal bases from slicing]\label{thm:bases}
The cascade's slicing structure, restricted to a 4D subspace, determines a family of
orthogonal bases for $\mathbb{R}^4$. Each choice of slicing axis within the 4D subspace
defines a measurement basis. Two measurements are complementary if and only if their slicing
axes are orthogonal in $\mathbb{R}^4$.
\end{theorem}
\begin{proof}
The slicing integral decomposes $B^{d+1}$ along a chosen axis $e$. Restricted to a 4D
subspace, the choice of $e$ selects a direction in $\mathbb{R}^4$. Different choices of $e$
give different decompositions of $\mathbb{R}^4$ into (axis)$\times$(3-ball). Two axes $e_1$,
$e_2$ are orthogonal if and only if the corresponding decompositions share no common axis,
which is the definition of complementary observables.
\end{proof}

\subsection{Non-commutativity from sequential slicing}

Slicing along axis $e_1$ then $e_2$ is not the same as slicing along $e_2$ then $e_1$,
because the first slice discards information about the first axis before the second slice
extracts information about the second.

\begin{theorem}[Commutator from slicing angle]\label{thm:commutator}
Let $e_1$, $e_2$ be two slicing axes in the 4D subspace with angle $\theta$ between them.
The order-dependence of sequential slicing vanishes if and only if $\theta=0$ or $\theta=\pi$.
For $\theta=\pi/2$, the order-dependence is maximal.
\end{theorem}

\begin{remark}
This recovers the structure of the canonical commutation relation $[x,p]=i\hbar$. The
position and momentum operators correspond to orthogonal slicing axes ($\theta=\pi/2$), and
the factor $i$ arises from complexification (Section~\ref{sec:complex}). The cascade does
not produce $\hbar$ directly; $\hbar$ enters as a scale factor when the geometric rates are
matched to physical units (Section~\ref{sec:hbar}).
\end{remark}

\section{The Born Rule from Concentration of Measure}

\subsection{Concentration on the high-dimensional sphere}

For a 1-Lipschitz function $f$ on $S^{d-1}(\sqrt{d})$:
\[
\Pr\big[|f-\mathbb{E}f|>t\big]\leq 2\exp(-t^2/2).
\]
This exponential concentration is the engine behind the Born rule.

\subsection{Deriving the Born rule}

\begin{theorem}[Born rule from cascade geometry]\label{thm:born}
Let $|\psi\rangle$ be a state on $S^{d-1}$ and let $P_k$ be the orthogonal projection onto
a $k$-dimensional subspace. The fraction of the uniform measure on $S^{d-1}$ that lies
within distance $\varepsilon$ of the subspace defined by $P_k$ converges, as $d\to\infty$,
to $|\langle\psi|P_k|\psi\rangle|^2=\|P_k|\psi\rangle\|^2$. This is the Born rule.
\end{theorem}
\begin{proof}
The proof follows the measure-concentration approach to Gleason's theorem~\cite{fuchs,clifton}.
On $S^{d-1}$, the uniform measure assigns to any great-subsphere $S^{k-1}$ a tubular
neighbourhood whose volume fraction is determined by the angle between the state vector and
the subspace. By concentration of measure, this fraction is sharply peaked at $\cos^2\theta$
where $\theta$ is the angle between $|\psi\rangle$ and the nearest point in the subspace.
\end{proof}

\subsection{What the cascade adds to Gleason}

Gleason's theorem~\cite{gleason} proves that the Born rule is the unique probability measure
on the lattice of subspaces of a Hilbert space of dimension $\geq 3$, but it assumes the
Hilbert space framework. The cascade provides the antecedent: the reason we have a Hilbert
space at all (Section~2), the reason for Gaussian measure (Section~3), and the reason for
orthogonal decomposition (Section~4).

\section{Complexification: The Forced Precession}
\label{sec:complex}

The cascade as described in~\cite{paperI} is real: it operates on $\mathbb{R}^d$. Quantum
mechanics requires complex amplitudes. This section shows how complex structure emerges from
the cascade, and---crucially---why the precession angle is not a free parameter.

\subsection{The forced precession}

\begin{theorem}[Forced precession]\label{thm:precession}
The angle between consecutive slicing axes is $\alpha = \pi/2$, forced by the cascade's
orthogonality axiom. No free parameter enters.
\end{theorem}
\begin{proof}
The cascade's orthogonality axiom is: each new slicing direction is perpendicular to all
previously integrated directions.

Applied to the slicing integral at step $d$: the slicing axis $e_{d+1}$ is perpendicular to
the equatorial hyperplane, forcing the half-integer argument in
$B(\frac{1}{2},\cdot)$, which gives $\Gamma(\frac{1}{2})=\sqrt{\pi}$ in the recurrence
(Theorem~3.1 of~\cite{paperI}).

Applied to consecutive axes: $e_{k+1}\perp e_k$ for all $k$ (the new axis is perpendicular
to the direction just integrated out), giving
$\alpha=\angle(e_k,e_{k+1})=\pi/2$.

Both conclusions follow from the
single axiom ``new slicing directions are perpendicular to all previously integrated
directions,'' applied at steps $d+1$ and $d$ respectively. The cascade's axiom forces both
simultaneously; there is no independent second assumption.

More precisely: the axiom is $e_k\perp e_j$ for all $k<j$ (each new slicing direction is
perpendicular to all earlier ones, which have been integrated out). For consecutive pairs
$(k,k+1)$: $\alpha=\angle(e_k,e_{k+1})=\pi/2$.
\end{proof}

\begin{corollary}[No free parameter in the precessing cascade]\label{cor:nofree}
The precessing cascade has no free geometric parameters. The precession angle $\alpha=\pi/2$
is determined by the cascade's orthogonality axiom alone.
\end{corollary}

\begin{remark}[Relationship to the $\sqrt{\pi}$ derivation]
Theorem~\ref{thm:precession} does not derive $\alpha=\pi/2$ from $\sqrt{\pi}$ or vice
versa. Both are consequences of the orthogonality axiom. The relationship is: the same
axiom ($e_k\perp$ all previously integrated directions) forces $\Gamma(\tfrac{1}{2})
=\sqrt{\pi}$ when applied to the integral over a single slicing direction, and forces
$\alpha=\pi/2$ when applied to the angle between consecutive slicing directions. They are
two theorems with a common hypothesis, not a chain where one implies the other.
\end{remark}

\subsection{Two quarter-turns yield the imaginary unit}

\begin{theorem}[Complex structure from forced precession]\label{thm:complex}
Let the cascade precess with $\alpha=\pi/2$ (forced by Theorem~\ref{thm:precession}). Then
the composition of two consecutive slicing operations, restricted to the plane of precession,
acts as multiplication by $-1$. The cascade's state space is naturally identified with a
complex Hilbert space, where the precession plane defines the complex structure $J$
($J^2=-\mathrm{Id}$).
\end{theorem}
\begin{proof}
A quarter-turn about $e_1$ maps $e_2\mapsto -e_1$ in the $e_1$-$e_2$ plane. A quarter-turn
about $e_2$ then maps $e_1\mapsto e_2$. The composition is a half-turn, i.e., multiplication
by $-1$. Define $J:e_1\mapsto e_2$, $e_2\mapsto -e_1$. Then $J^2=-\mathrm{Id}$, the
algebraic definition of a complex structure on $\mathbb{R}^2$. Extending to all pairs of
precessing axes partitions $\mathbb{R}^{2n}$ into $n$ complex dimensions $\mathbb{C}^n$,
with the cascade's quarter-turn providing the geometric origin of $i=e^{i\pi/2}$.

Since $\alpha=\pi/2$ is forced (Theorem~\ref{thm:precession}), this complex structure is
not introduced as an assumption but derived.
\end{proof}

\begin{remark}[The propagator phase at each cascade level]\label{rem:bott-phases}
The forced precession of $\pi/2$ per step accumulates a total propagator phase of
$(d-4)\cdot\pi/2$ from the observer at $d=4$ to level $d$. The phase at each level is
therefore determined entirely by $d\bmod 8$:
\begin{center}
\begin{tabular}{ccc}
\toprule
$d\bmod 8$ & Spinor type & Phase \\
\midrule
4 & Weyl (complex) & $e^{i\cdot 0} = +1$ \\
5 & Dirac (complex) & $e^{i\pi/2} = i$ \\
6 & Weyl (complex) & $e^{i\pi} = -1$ \\
7 & Majorana (real) & $e^{i3\pi/2} = -i$ \\
0 & Majorana (real) & $e^{i2\pi} = +1$ \\
1 & Majorana (real) & $e^{i5\pi/2} = i$ \\
2 & Majorana (real) & $e^{i3\pi} = -1$ \\
3 & Majorana (real) & $e^{i7\pi/2} = -i$ \\
\bottomrule
\end{tabular}
\end{center}
The propagator phase is a topological invariant: it depends only on $d\bmod 8$, not on the
metric. The phase column follows from Theorem~\ref{thm:complex} alone: it is the accumulated
precession $(d-4)\cdot\pi/2$ reduced mod $2\pi$, with no external input. The spinor type
column follows from the Clifford algebra classification of $\mathrm{Cl}(1,d-1)$ under Bott
periodicity, an established result of algebraic topology.

The two columns make distinct claims. The Clifford classification establishes that layers
with $d\bmod 8\in\{4,5,6\}$ carry irreducibly complex spinors, while Majorana layers
($d\bmod 8\in\{0,1,2,3,7\}$) are real. This applies uniformly to all three complex residues.
The phase column then distinguishes within the complex layers: $d\bmod 8=4$ (Weyl) has
phase $+1$, meaning the propagator returns to the real axis after one Bott period;
$d\bmod 8\in\{5,6\}$ (Dirac and Weyl) have phases $i$ and $-1$, meaning the propagator
does not return to the real axis within one period. The distinction between $\{4\}$ and
$\{5,6\}$ is therefore not spinor irreducibility---all three residues share that
property---but propagator phase non-triviality, which is derived from
Theorem~\ref{thm:complex} alone.
\end{remark}

\subsection{Interference}

Complex structure immediately gives interference. If two cascade paths arrive at the same
projected state with different accumulated precession angles $\varphi_1$ and $\varphi_2$,
their contributions are $e^{i\varphi_1}$ and $e^{i\varphi_2}$. The 4D observer measures:
\[
|e^{i\varphi_1}+e^{i\varphi_2}|^2 = 2+2\cos(\varphi_1-\varphi_2),
\]
exhibiting constructive interference when $\varphi_1=\varphi_2$ and destructive when
$\varphi_1-\varphi_2=\pi$.

\section{Time as Orthogonal Descent}

\subsection{The identification}

We identify the cascade's slicing direction with time. Each step of the cascade---integrating
out one orthogonal direction---corresponds to one unit of temporal evolution. The cascade
does not happen in time; the cascade is time. The arrow of time is the direction of descent,
which is irreversible: the slicing integral loses information about the integrated-out
direction.

\subsection{The lapse function}

The cascade's slicing recurrence gives the ratio of adjacent volumes:
\[
V_{d+1}/V_d = \sqrt{\pi}\cdot R(d),\qquad R(d) = \frac{\Gamma((d+1)/2)}{\Gamma((d+2)/2)}.
\]
This is the cascade's lapse function: $N(d)=\sqrt{\pi}\cdot R(d)\approx\sqrt{2\pi/d}$ in
the Stirling regime. Properties: $N(d)\to 0$ as $d\to\infty$ (time does not exist at the
cascade's starting point); $N(d)$ monotonically increasing as $d$ decreases; $N(4)=3\pi/8
\approx 1.178$ (the 4D observer lives in a moderate regime; the Stirling approximation
$\sqrt{\pi/2}\approx 1.25$ is 6\% high).

\subsection{The discrete propagator}

\begin{theorem}[Cascade propagator]\label{thm:propagator}
The composition of cascade steps from $d=D$ down to $d=d'$ defines a propagator
$K(D,d')=\prod_{j=d'}^{D-1}L(j)$. In the precessing cascade, each $L(j)$ acquires the
forced phase factor $e^{i\alpha}=e^{i\pi/2}=i$, giving
$K(D,d')=\prod_j|L(j)|\cdot i^{D-d'}$.
\end{theorem}

\subsection{Continuum limit: the Schr\"{o}dinger equation}

In the continuum limit where $d$ is treated as a continuous parameter and the precessing
cascade has a smooth phase, the propagator equation becomes:
\[
i\frac{d\psi}{dt} = H(t)\,\psi(t),
\]
where $t$ parametrises the cascade descent, $\psi$ is the projected state in the 4D
subspace, and $H(t)$ encodes the local slicing rate and precession. The factor $i$ arises
from Theorem~\ref{thm:complex}. This is the Schr\"{o}dinger equation.

\begin{remark}
The arrow of time is structural, not thermodynamic. The cascade descends because the slicing
integral compresses $(d+1)\to d$; no inverse operation restores the integrated-out direction
without supplying external information.
\end{remark}

\section{Compactification and the Geometry of Temporal Descent}
\label{sec:compactification}

Section~7 identifies time with slicing and states that the arrow of time arises from the
irreversibility of integrating out a direction. This section gives that identification
concrete geometric content: each slicing step does not eliminate a direction but
asymptotically suppresses it, with an exact compactification radius determined by the Beta
function.

\subsection{The effective radius of each temporal step}

\begin{theorem}[Compactification radius]\label{thm:reff}
At each slicing step $d$, the integrated-out direction retains an effective radius
\[
R_{\rm eff}(d) = \frac{1}{\sqrt{d+3}}.
\]
\end{theorem}
\begin{proof}
The effective radius is the RMS extent of the integrated-out coordinate under the
slicing measure:
\[
R_{\rm eff}^2 = \langle x^2\rangle
= \frac{\int_{-1}^{1} x^2(1-x^2)^{d/2}\,dx}{\int_{-1}^{1}(1-x^2)^{d/2}\,dx}.
\]
Substituting $t = x^2$, the numerator is $B(3/2,\,d/2+1)$ and the denominator is
$B(1/2,\,d/2+1)$. Their ratio is
\[
\langle x^2\rangle = \frac{B(3/2,\,d/2+1)}{B(1/2,\,d/2+1)}
= \frac{\Gamma(3/2)}{\Gamma(1/2)}\cdot\frac{\Gamma(d/2+3/2)}{\Gamma(d/2+5/2)}
= \frac{1}{2}\cdot\frac{1}{d/2+3/2}
= \frac{1}{d+3},
\]
where the first factor is $\Gamma(3/2)/\Gamma(1/2) = 1/2$ and the second uses
$\Gamma(x+1) = x\,\Gamma(x)$ once.
Therefore $R_{\rm eff} = 1/\sqrt{d+3} \to 0$ as $d\to\infty$.
\end{proof}

The compactification is asymptotic: $R_{\rm eff}\neq 0$ at any finite $d$. The weight
function $(1-x^2)^{d/2}$ is smooth and positive on $(-1,1)$ for all finite $d$. There is
no sharp boundary, no horizon, no topology change---only exponential suppression. The word
``asymptotic'' is essential: the compactification limit is approached but never reached at
any finite step.

\begin{center}
\begin{tabular}{cccc}
\toprule
$d$ & $R_{\rm eff}$ & $N(d)$ & Physical role \\
\midrule
4   & 0.37796 & 1.17810 & Observer dimension \\
5   & 0.35355 & 1.06667 & $S^3$ boundary layer \\
7   & 0.31623 & 0.91429 & Area maximum $d_0$ \\
19  & 0.21320 & 0.56755 & First threshold $d_1$ \\
217 & 0.06742 & 0.16997 & Second threshold $d_2$ \\
\bottomrule
\end{tabular}
\end{center}

\subsection{The scale coincidence}

The cascade has two natural length scales at each step: the compactification radius
$R_{\rm eff}(d) = 1/\sqrt{d+3}$ from the Beta function, and the Gaussian width
$\sigma(d) = 1/\sqrt{d}$ of the slicing integrand (Lemma~\ref{lem:gaussian}). These
are defined by different objects---one by the second moment, one by the integrand's
curvature at $x = 0$---yet they coincide asymptotically.

\begin{theorem}[Scale coincidence]\label{thm:scale}
\[
\frac{R_{\rm eff}(d)}{\sigma(d)} = \sqrt{\frac{d}{d+3}} \to 1 \quad\text{as }
d\to\infty,
\]
with correction $O(1/d)$: $R_{\rm eff}/\sigma = 1 - \frac{3}{2d} + O(d^{-2})$.
\end{theorem}
\begin{proof}
Direct computation from Lemma~\ref{lem:gaussian} and Theorem~\ref{thm:reff}.
\end{proof}

\begin{corollary}[Second universal cascade constant]\label{cor:erf}
The fraction of the slicing integrand's content within the compactification radius
satisfies
\[
F(d) = \frac{\int_{|x|<R_{\rm eff}(d)}(1-x^2)^{d/2}\,dx}
{\int_{-1}^{1}(1-x^2)^{d/2}\,dx}
\to \operatorname{erf}\!\left(\frac{1}{\sqrt{2}}\right) = 0.68269\ldots
\quad\text{as }d\to\infty,
\]
a second universal constant of the cascade alongside $\sqrt{\pi}$, determined entirely
by the coincidence $R_{\rm eff}/\sigma\to 1$.
\end{corollary}
\begin{proof}
In the Gaussian limit (Lemma~\ref{lem:gaussian}), the integrand is
$e^{-dx^2/2}$ with width $\sigma = 1/\sqrt{d}$. The fraction within
$R_{\rm eff} = 1/\sqrt{d+3}$ is
$\operatorname{erf}(R_{\rm eff}/(\sigma\sqrt{2}))$. As $d\to\infty$,
$R_{\rm eff}/(\sigma\sqrt{2})\to 1/\sqrt{2}$, giving
$\operatorname{erf}(1/\sqrt{2})$.
\end{proof}

\begin{remark}[The cascade's effective geometric origin]
After $d_2 = 217$, sphere areas satisfy $\Omega_d < 10^{-120}$---geometrically
indistinguishable from zero. The cascade's second threshold $d_2 = 217$ and its formal
starting point $d = \infty$ are therefore the same geometric point: the cascade begins at
geometric nothing, and nothing remains after $d_2$. The limit $d\to\infty$ in
Corollary~\ref{cor:erf} is consequently exact, not merely asymptotic, when applied to the
cascade tower: there is no geometric content beyond $d_2 = 217$ to distinguish the two
points. The constant $\operatorname{erf}(1/\sqrt{2})$ is a property of the cascade's
origin, evaluated at its effective terminus.
\end{remark}

\begin{remark}[Structural foreshadowing]
The cascade produces exactly two universal constants from its geometry:
$\sqrt{\pi} = \Gamma(1/2)$ from orthogonality, and $\operatorname{erf}(1/\sqrt{2}) =
0.68269$ from the scale coincidence. The first generates the cosmological constant.
The second is within $0.15\%$ of $(\pi-1)/\pi = 0.68169$, which Part~V derives as the
dark energy fraction $\Omega_\Lambda$. Both constants are properties of the
Gamma function evaluated at its cascade-distinguished points; neither is fitted.
\end{remark}

\subsection{Irreversibility as asymptotic suppression}

The identification of time with slicing (Section~7.1) asserts that the slicing integral
loses information about the integrated-out direction. Theorem~\ref{thm:reff} makes this
precise: the direction is not lost but suppressed to scale $R_{\rm eff}=1/\sqrt{d+3}$.
From the perspective of the $(d+1)$-dimensional geometry, the direction's contribution is
exponentially localised near $x=0$ with Gaussian width $\sigma\approx 1/\sqrt{d}$
(Lemma~\ref{lem:gaussian}). From the perspective of the $d$-dimensional residual geometry,
the direction has been compactified.

The arrow of time acquires a quantitative meaning: each step of temporal descent reduces one
direction's effective extent by a factor that depends on $d$. Early in the cascade (large
$d$), $R_{\rm eff}$ is small and the suppression per step is mild. Late in the cascade
(small $d$), $R_{\rm eff}$ is larger and each step represents a more substantial geometric
change. At the observer's dimension $d=4$, $R_{\rm eff}=1/\sqrt{7}\approx 0.378$: each
temporal step compactifies a direction to roughly 38\% of the unit ball's radius.

The irreversibility of slicing is now geometrically transparent. To reverse a slicing step
would require decompactifying a direction from $R_{\rm eff}$ back to unit radius. This
requires external information: the specific correlations between the integrated-out
coordinate and the remaining geometry, which are asymptotically suppressed in the slicing
integral. The arrow of time is the asymptotic approach to the compactification limit,
viewed from the descending side.

\subsection{Unitarity and information}

The cascade dynamics are unitary: the Schr\"{o}dinger equation (Section~7.4) preserves
norms. Yet each slicing step appears irreversible (Section~7.1). These two statements are
compatible because they describe different levels.

The full cascade geometry from $d=D$ down to $d=d'$ is described by a unitary propagator
$K(D,d')$ (Theorem~\ref{thm:propagator}). No information is lost in the total evolution.
A single slicing step, viewed from the $d$-dimensional residual geometry, appears to lose
information because the observer cannot access the compactified direction at scale
$R_{\rm eff}=1/\sqrt{d+3}$. The lost information is not destroyed; it is encoded in the
correlations between the compactified direction and the boundary.

This is made precise by boundary dominance: $\Omega_{d-1}/V_d=d$ (Theorem~3.1
of~\cite{paperI}). The boundary $S^{d-1}$ carries a fraction $d/(d+1)$ of the content at
each layer. Information that appears lost from the interior is retained on the boundary.
The cascade is unitary as a whole; it looks irreversible from any single layer because the
boundary of that layer encodes information the interior observer cannot directly access.

\begin{remark}[Comparison with the measurement problem]
The structure is suggestive: unitary global evolution, apparent irreversibility from a
restricted perspective, information retained on the boundary. These are the ingredients of
both the quantum measurement problem and the black hole information paradox. The cascade
provides them from geometry alone. Whether this structural parallel constitutes a resolution
of either problem is addressed in Section~\ref{sec:open}.
\end{remark}

\section{The Role of $\hbar$}
\label{sec:hbar}

\subsection{$\hbar$ as the unit-matching constant}

By Theorem~\ref{thm:precession}, the precession angle is fixed: $\alpha=\pi/2$. The cascade
is purely geometric and produces dimensionless quantities. Physical quantum mechanics has a
dimensionful constant $\hbar$ with units of action. It enters as:
\[
\hbar = \frac{\alpha}{\Delta t} = \frac{\pi/2}{\Delta t},
\]
where $\Delta t$ is the physical duration assigned to one cascade step. Since $\alpha$ is
now fixed by geometry, $\hbar$ encodes only the identification of cascade steps with
physical time increments. It is a unit-matching constant, not a free geometric parameter.

\subsection{The classical limit}

The classical limit $\hbar\to 0$ corresponds to $\Delta t\to\infty$ at fixed
$\alpha=\pi/2$: the observer's temporal resolution is too coarse to track the precession.
Phases decorrelate, and the 4D observer sees incoherent (classical) projections. This is
geometric decoherence.

\section{Tensor Products and Entanglement}

\subsection{Tensor structure from iterated slicing}

The slicing recurrence decomposes $V_{d+1}=V_d\times(\text{fibre})$, where the fibre is
the one-dimensional integral over the perpendicular coordinate. Iterating $k$ times:
$V_{d+k}=V_d\times\text{fibre}_1\times\cdots\times\text{fibre}_k$. This is a tensor
product: the $(d+k)$-dimensional space factors into a $d$-dimensional base and $k$
one-dimensional fibres.

\subsection{Entanglement}

\begin{theorem}[Generic entanglement]\label{thm:entangle}
For $d\geq 2$, the cascade state after one slicing step is generically entangled between
the base and fibre subsystems. Separable states have measure zero in the cascade's state
space.
\end{theorem}
\begin{proof}
The slicing cross-section at height $x$ has radius $\sqrt{1-x^2}$, creating a correlation
between the fibre coordinate $x$ and the base radius. A product state would require the
cross-sectional radius to be independent of $x$, which occurs only for $d=0$. For $d\geq 2$,
the coupling term $(1-x^2)^{d/2}$ is non-trivial, and the set of separable states is a
lower-dimensional submanifold.
\end{proof}

The compactification interpretation (Section~\ref{sec:compactification}) sharpens the
entanglement structure. The fibre coordinate $x$ is compactified to effective radius
$R_{\rm eff}=1/\sqrt{d+3}$, but the base radius $\sqrt{1-x^2}$ depends on $x$ throughout
the compactification region. The base--fibre entanglement is generated by this $x$-dependent
coupling and persists at all scales above $R_{\rm eff}$. Below $R_{\rm eff}$, the fibre is
effectively frozen and the entanglement is inaccessible to a $d$-dimensional observer---but
it is retained in the boundary correlations.

\subsection{Bell correlations}

The 4D observer, seeing projections of a $d=217$ geometry, observes correlations between
subsystems that cannot be explained by any local hidden-variable model---because the
correlations arise from the high-dimensional geometry that the observer cannot directly
access. The hidden geometry is not local (it extends through 217 dimensions) and not
classical (it has the complex structure of Theorem~\ref{thm:complex}).

\section{Summary: From Cascade to Quantum Mechanics}

\begin{center}
\begin{tabular}{p{5.5cm}p{5cm}c}
\toprule
QM Postulate & Cascade Origin & Section \\
\midrule
State space is a Hilbert space & Unit sphere from cascade geometry & 2 \\
States are unit vectors & Normalisation from ball boundary & 2 \\
Arena is the sphere, not the ball & Boundary dominance $\Omega/V=d$ & 2.3 \\
Gaussian wavefunctions & Projection from high dimensions (CLT) & 3 \\
Gaussian integrand $(1-x^2)^{d/2}$ & Concentration of slicing measure & 3.1 \\
Normalisation by $\sqrt{\pi}$ & Orthogonal compression constant & 3 \\
Orthogonal measurement bases & Slicing axis choices & 4 \\
Measurement as suppression & $R_{\rm eff}=1/\sqrt{d+3}$ & 4.1, 8 \\
Non-commuting observables & Sequential slicing order-dependence & 4 \\
Born rule $p=|\langle\phi|\psi\rangle|^2$ & Concentration of measure on sphere & 5 \\
Complex amplitudes & Forced precession $\alpha=\pi/2\Rightarrow J^2=-\mathrm{Id}$ & 6 \\
No free geometric parameters & $\alpha=\pi/2$ forced; same axiom as $\sqrt{\pi}$ & 6.1 \\
Phase classification by $d\bmod 8$ & Forced precession (phase); Bott periodicity (spinor) & 6.5 \\
Interference & Phase accumulation from forced precession & 6.3 \\
Unitary time evolution & Cascade propagator (discrete) & 7 \\
Schr\"{o}dinger equation & Continuum limit of propagator & 7.4 \\
Arrow of time & Asymptotic compactification per step & 7, 8 \\
Scale coincidence $R_{\rm eff}/\sigma\to 1$ & Compactification + Gaussian width & 8.2 \\
Second cascade constant $\operatorname{erf}(1/\sqrt{2})$ & Scale coincidence at $d\to\infty$ & 8.2 \\
Irreversibility is asymptotic & $R_{\rm eff}>0$; no sharp loss & 8.3 \\
Unitarity + apparent irreversibility & Global unitary; boundary retains info & 8.4 \\
$\hbar$ as quantum scale & Unit matching: $\hbar=(\pi/2)/\Delta t$ & 9 \\
Classical limit & Coarse time resolution $\Rightarrow$ no phase & 9.2 \\
Tensor product structure & Iterated slicing decomposition & 10 \\
Entanglement & Non-factorisable cascade states & 10.2 \\
Bell inequality violation & Non-local high-dimensional geometry & 10.3 \\
\bottomrule
\end{tabular}
\end{center}

Every row is derived from two inputs: (1) the cascade geometry of~\cite{paperI}, and
(2) the assumption that the observer has access to 4 dimensions. The compactification
results (Section~\ref{sec:compactification}) are proved in this paper from the Beta
function, building on the boundary dominance identity of~\cite{paperI}, Theorem~3.1.
The phase classification row (6.5) is derived from the forced precession of this paper
combined with Bott periodicity.

\section{What This Paper Does and Does Not Do}

\textbf{Does:}
\begin{itemize}
\item Shows that a 4D observer in the cascade geometry recovers the structural framework of
quantum mechanics.
\item Proves that the precession angle $\alpha=\pi/2$ is forced by the cascade's
orthogonality axiom, eliminating all free geometric parameters
(Theorem~\ref{thm:precession}). The proof explicitly distinguishes the two applications of
the orthogonality axiom (one forcing $\sqrt{\pi}$ in the integral, one forcing $\alpha$)
and shows they are sequential applications of the same axiom, not a circular argument.
\item Derives the Gaussian character of wavefunctions from projection, with the elementary
mechanism identified: the slicing integrand $(1-x^2)^{d/2}$ is itself approximately Gaussian
(Lemma~\ref{lem:gaussian}).
\item Provides a geometric origin for complex amplitudes via forced axis precession.
\item Identifies time with the cascade's slicing direction, yielding the arrow of time and
a natural lapse function.
\item Recovers the Schr\"{o}dinger equation as the continuum limit of the cascade propagator.
\item Proves the compactification radius $R_{\rm eff}=1/\sqrt{d+3}$ from the Beta function
(Theorem~\ref{thm:reff}) and develops its physical interpretation: each temporal step is an
asymptotic compactification. The arrow of time, the nature of measurement, and the
compatibility of unitarity with apparent irreversibility all receive concrete geometric
content from this identification.
\item Establishes the scale coincidence $R_{\rm eff}/\sigma\to 1$ (Theorem~\ref{thm:scale})
and derives the second universal cascade constant $\operatorname{erf}(1/\sqrt{2}) = 0.68269$
(Corollary~\ref{cor:erf}).
\item Shows that boundary dominance ($\Omega_{d-1}/V_d=d$) justifies the primacy of the
unit sphere as the observer's state space.
\item Derives the propagator phase classification by $d\bmod 8$ from the forced precession
and Bott periodicity (Remark~\ref{rem:bott-phases}).
\end{itemize}

\textbf{Does not:}
\begin{itemize}
\item Derive why we observe 4 dimensions. This is an empirical input.
\item Derive the specific Hamiltonian of any physical system. The cascade gives the
framework but not the content.
\item Derive the value of $\hbar$. The geometric factor $\pi/2$ is derived; the physical
time increment $\Delta t$ is not.
\item Derive quantum field theory. The cascade gives single-particle QM; extension to fields
requires additional structure.
\item Use any result from physics. The only inputs are the cascade's geometry and the
observer's dimensionality.
\end{itemize}

\section{Open Questions}
\label{sec:open}

\begin{enumerate}
\item \textbf{Quantum field theory.} Whether the dimensions between successive thresholds
encode specific field content is open.
\item \textbf{Physical time increment.} The geometric value $\alpha=\pi/2$ is derived. The
remaining question is how the physical time increment $\Delta t$ is fixed, which would
determine $\hbar$.
\item \textbf{Lorentzian signature.} The cascade is real and Euclidean. The forced precession
$\alpha=\pi/2$ produces $J^2=-\mathrm{Id}$. Whether the identification $g_{tt}=(iN)^2=-N^2$
can be derived as a theorem is open.
\item \textbf{The measurement problem.} The cascade provides a natural account of measurement
as asymptotic compactification: the measured direction is suppressed to $R_{\rm eff}=
1/\sqrt{d+3}$, not sharply projected. Whether this smooth suppression---combined with the
boundary's retention of information ($\Omega/V=d$)---replaces the projection postulate is
open. The structural ingredients (global unitarity, apparent local irreversibility, boundary
encoding) are suggestive but not yet a derivation.
\item \textbf{Decoherence from compactification.} The compactification radius $R_{\rm eff}$
defines a natural decoherence scale: correlations at distances below $R_{\rm eff}$ are
inaccessible to the lower-dimensional observer. Whether this geometric decoherence reproduces
the quantitative predictions of environment-induced decoherence~\cite{zurek} is open.
\end{enumerate}

\begin{thebibliography}{9}
\bibitem{paperI} RTAC, \emph{The Cascade Series --- Part 0: Scale Variance from
Orthogonality: How the Unit Ball Generates $10^{120}$ Orders of Magnitude}, March 2026.
\bibitem{ledoux} M.~Ledoux, \emph{The Concentration of Measure Phenomenon}, AMS
Mathematical Surveys and Monographs 89, 2001.
\bibitem{milman} V.~D.~Milman and G.~Schechtman, \emph{Asymptotic Theory of
Finite-Dimensional Normed Spaces}, Lecture Notes in Mathematics 1200, Springer, 1986.
\bibitem{fuchs} C.~A.~Fuchs, ``Quantum Mechanics as Quantum Information (and only a little
more),'' arXiv:quant-ph/0205039, 2002.
\bibitem{clifton} R.~Clifton, J.~Bub, and H.~Halvorson, ``Characterizing quantum theory in
terms of information-theoretic constraints,'' \emph{Foundations of Physics} \textbf{33},
1561--1591 (2003).
\bibitem{gleason} A.~M.~Gleason, ``Measures on the closed subspaces of a Hilbert space,''
\emph{Journal of Mathematics and Mechanics} \textbf{6}, 885--893 (1957).
\bibitem{zurek} W.~H.~Zurek, ``Quantum Darwinism,'' \emph{Nature Physics} \textbf{5},
181--188 (2009).
\bibitem{hardy} L.~Hardy, ``Quantum theory from five reasonable axioms,''
arXiv:quant-ph/0101012, 2001.
\bibitem{adm} R.~Arnowitt, S.~Deser, and C.~W.~Misner, ``The dynamics of general
relativity,'' in \emph{Gravitation: An Introduction to Current Research}, Wiley, 1962.
\end{thebibliography}

\end{document}